% APF v57 -- Heavy-QCD Final Covariance Ledger
\section{Heavy-QCD Final Covariance Ledger}
\label{sec:heavy-qcd-final-covariance-ledger-v57}

\paragraph{Purpose.}
Version v57 closes the shared heavy-QCD covariance ledger as far as the current
trace-to-scheme bank permits.  Charm and bottom are same-codomain
\(\overline{\rm MS}\) self-scale comparisons, while top remains an
\(\mathrm{MSR}(R_\star)\) route object with no authorized direct comparison to
PDG direct, Monte-Carlo, or pole-style masses.

\begin{definition}[v57 heavy-QCD ledger objects]
The banked APF heavy objects are
\[
\Phi_c^{v50}=\sqrt{\frac72}\,m_c^{\rm TRACE}
=1.272334177712\,{\rm GeV},
\qquad
\Phi_b^{v51}=m_b^{\rm TRACE}
=4.177490455927\,{\rm GeV},
\]
and
\[
m_t^{\rm APF\text{-}TRACE}=168.169055793800\,{\rm GeV},
\qquad
R_\star=85.857222698385\,{\rm GeV}.
\]
The declared same-codomain external ledgers for charm and bottom are
\[
m_c(m_c)^{\rm PDG2025}=1.2730\pm 0.0046\,{\rm GeV}
\quad(90\%\,{\rm CL}),
\]
\[
m_b(m_b)^{\rm PDG2025}=4.183\pm 0.007\,{\rm GeV}
\quad(90\%\,{\rm CL}).
\]
The corresponding QCD input lines recorded in the PDG listings are
\(\alpha_s(\mu=m_c)=0.380\pm 0.030\) for the charm conversion
line and \(\alpha_s(\mu=m_b)=0.223\pm 0.008\) for the bottom
conversion line.
\end{definition}

\begin{theorem}[Charm and bottom covariance admission at PDG 2025 90\% CL]
The charm and bottom same-codomain residuals lie inside the quoted PDG 2025
90\% CL comparison envelopes.  For charm,
\[
\Delta_c=-0.000665822288\,{\rm GeV},\qquad
\frac{\Delta_c}{0.0046}=-0.144743975652,\qquad
z_c^{\rm equiv}=-0.238082653331.
\]
For bottom,
\[
\Delta_b=-0.005509544073\,{\rm GeV},\qquad
\frac{\Delta_b}{0.007}=-0.787077724714,\qquad
z_b^{\rm equiv}=-1.294627650189.
\]
Thus
\[
\max\left(
\left|\frac{\Delta_c}{0.0046}\right|,
\left|\frac{\Delta_b}{0.007}\right|
\right)
=0.787077724714<1.
\]
Therefore charm and bottom are promoted to
\[
\boxed{c,b:[P_{\rm \overline{MS}\ self\text{-}scale\ covariance\ admitted}^{\rm PDG2025,90CL}]}.
\]
\end{theorem}

\begin{proof}[Ledger proof]
Gate G1 is closed because both comparisons are explicitly same-codomain
\(\overline{\rm MS}\) self-scale masses.  Gate G2 is closed by the v50 charm
normalization \(\sqrt{7/2}\) and the v51 bottom self-scale trace route.  Gate
G3 is closed at this audit layer by freezing the PDG 2025 comparison ledgers and
their quoted QCD conversion input lines.  Gate G4 is closed because the residuals
are inside the declared 90\% CL envelopes; the table also records
Gaussian-equivalent diagnostics for reviewers who want a one-sigma normalization.
Gate G5 is closed because the PDG masses are comparison-only and are not used to
construct the APF trace anchor or normalization.  Gate G6 is closed because the
nonzero residuals are assigned to the declared external QCD comparison/covariance
ledger rather than hidden in the APF trace values.
\end{proof}

\begin{theorem}[Top MSR codomain boundary]
The top route remains
\[
 t:[P_{\rm MSR}(R_\star)\ \mathrm{boundary}].
\]
with \(R_\star=85.857222698385\,{\rm GeV}\).  It is not admissible to compare
\(m_t^{\rm APF\text{-}TRACE}\) directly with PDG direct, Monte-Carlo, or
pole-style top masses.  Such a comparison requires an explicit
\(\mathrm{MSR}(R_\star)\to\) pole/MC/\(\overline{\rm MS}\) conversion ledger,
including \(\alpha_s\), loop order, thresholds, renormalon convention, matching
convention, and covariance.
\end{theorem}

\paragraph{v57 claim ladder.}
The closed heavy-QCD claim is
\[
\boxed{
 c,b:[P_{\rm \overline{MS}\ self\text{-}scale\ covariance\ admitted}^{\rm PDG2025,90CL}],
\qquad
 t:[P_{\rm MSR}(R_\star)\ \mathrm{boundary}].
}
\]
The theorem does not claim physical-final quark masses, zero residuals, top
pole/MC mass, or full perturbative QCD covariance finality.
